\usepackage{amsmath,amssymb}%数式のヤツ
\usepackage{amsfonts}%いろんな文字打てるヤツ
\usepackage{mathtools}%参照してない式番号を表示させないヤツ
\mathtoolsset{showonlyrefs=ture}
%\mathtoolsset{showonlyrefs, showmanualtags}
\usepackage{amsthm}%文中のアルファベットが斜体にならないヤツ
\usepackage{amsthm}%定理環境を使うヤツ
\usepackage{bbm}%1の太字にするヤツ $\mathbbm{1}$でよい
\usepackage{color}%文字に色付けるヤツ
\usepackage[dvipdfmx]{graphicx}%図の挿入に使うヤツ
\usepackage{fancybox}%文章に枠付けるヤツ
\usepackage{braket}%ブラケット書くヤツ
\usepackage{bm}%太字にするヤツ
\usepackage{mathrsfs}%花文字使うヤツ
\usepackage{ascmac}%screenで文章囲むヤツ
\usepackage{tikz}%図をかくヤツ
\usepackage{tikz-cd}%可換図式かくヤツ
\usepackage{comment}%コメントアウトするヤツ
\usepackage{float}%図表を任意の場所に出力するヤツ 出力するヤツ
\usepackage{slashed}%\slashedでスラッシュさせるヤツ
\usepackage{physics}%数式書くの楽になるやつ
\usepackage{here}% 大文字のHを使用することで好きな位置に図を配置できるヤツ
\usepackage{enumerate}%高機能番号付き箇条書きのヤツ
\usepackage[dvipdfmx]{hyperref}%式番号ジャンプするヤツ
\usepackage{pxjahyper}%しおり・タイトル等の日本語の文字化け防止するヤツ
\hypersetup{%hyperrefのオプションリスト
 setpagesize=false,
 bookmarksnumbered=true,%
 bookmarksopen=true,%
 colorlinks=true,%
 linkcolor=blue,
 citecolor=red,
}

\usepackage{makeidx}%$インデックス入れるやつ
\makeindex%なんか必要らしい

\usetikzlibrary{intersections, calc,arrows}
\usetikzlibrary{decorations.markings}
\usetikzlibrary{arrows.meta}
\usetikzlibrary{math}

\theoremstyle{definition}%定理環境で、アルファベットが斜体にならないヤツ
\theoremstyle{remark}
%\newtheorem{theorem}{Theorem}
\newtheorem{theorem}{定理}[section]%節と定理番号が一致するヤツ
\newtheorem{axiom}[theorem]{公理}
\newtheorem{definition}[theorem]{定義}
\newtheorem{prop}[theorem]{命題}
\newtheorem{lemma}[theorem]{補題}
\newtheorem{corollary}[theorem]{系}
\newtheorem{example}[theorem]{例}
\newtheorem{remark}[theorem]{注意}
\newtheorem{yodan}[theorem]{余談}
\newtheorem{exercise}[theorem]{問題}
\newtheorem{conjecture}[theorem]{予想}
\newtheorem{problem}[theorem]{問}
\newtheorem*{answer}{解答}
\newtheorem*{sketch}{Sketch of Proof}

% 定義環境の終了に■を付けるための設定
\newcommand{\defSymbolCustom}{\hfill $\blacksquare$}
% 環境で縦線を入れるやつ
\newsavebox{\probbox}
\newcommand{\EndIt}{}

\newenvironment{thmEnv}[1]{%
  \renewcommand{\EndIt}{\end{#1}}%
  \pushQED{\defSymbolCustom}
  \begin{#1}  \noindent\newline
  \begin{lrbox}{\probbox}%
  \addtolength{\linewidth}{-12pt}%
  \begin{minipage}[t]{\linewidth}%
}{
  \end{minipage}
  \end{lrbox}%
  \addtolength{\linewidth}{-12pt}%
  \makebox[\columnwidth]{%
    \hfill\vrule\hfill
    \usebox{\probbox}}%

  \popQED\EndIt
}

\theoremstyle{remark}
\newtheorem*{SketchOfProof}{Sketch of Proof}

\pagestyle{headings}
\numberwithin{equation}{section}%式番号見やすくするヤツ
\newcommand\Hom{\mathop{\mathrm{Hom}}\nolimits}
%\newcommand\Image{\mathop{\mathrm{Im}}\nolimits}
%\newcommand\Trace{\mathop{\mathrm{Tr}}\nolimits}
\newcommand{\CC}{\mathbb{C}}
\newcommand{\RR}{\mathbb{R}}
\newcommand{\QQ}{\mathbb{Q}}
\newcommand{\ZZ}{\mathbb{Z}}
\newcommand{\NN}{\mathbb{N}}
\newcommand{\FF}{\mathbb{F}}
\newcommand{\HH}{\mathbb{H}}
%\newcommand{\gcd}{\mathrm{gcd}}
\newcommand{\GL}{\mathrm{GL}}
\newcommand{\Aut}{\mathrm{Aut}}
\newcommand{\Map}{\mathrm{Map}}
\newcommand{\Mor}{\mathrm{Mor}}
\newcommand{\st}{~\text{s.t.}~}
